\documentclass[crop,border=7pt]{standalone}

\usepackage{graphicx}
\usepackage{float}
\usepackage{amsmath}       
\usepackage{scalefnt}
\usepackage[dvipsnames]{xcolor}
%% TIKZ
\usepackage{tikz}
\usepackage{pgfplots}
\pgfplotsset{compat=newest}
\usetikzlibrary{decorations.pathmorphing,decorations.text,decorations.pathreplacing}
\usetikzlibrary{patterns}
\usetikzlibrary{shapes,shapes.multipart,arrows,fit,calc,positioning,intersections,through}
\usetikzlibrary{backgrounds,fadings}

 
%%%%%%%%%%%%%%%%%%%%%%%%%%%%%%%%%%%%%%%%%%%%%%%%%%%%%%%%%%%%%%%%%%%%%%%%%%%%%%%%%%%%%%%%%%%%%%%%%%
\begin{document}
%  \tikzsetnextfilename{flow_chart}
 \tikzstyle{rect} = [draw=black, thick, rectangle, fill=yellow!30, text width=16em, text centered, minimum height=2em]
 \tikzstyle{elli} = [draw=black, thick, ellipse, fill= white!20, minimum height=2em, text width=10em]
\tikzstyle{circ} = [draw=black, thick, circle, fill=red!30, text width=3em, text centered]
\tikzstyle{diam} = [draw=black, thick, diamond, fill=green!30, text width=4.5em, text badly centered, inner sep=0pt]
\tikzstyle{trap} = [draw=black, thick, trapezium, trapezium left angle=70, trapezium right angle=110, fill=white!20, text width=6em, text centered, minimum height=2em]
\tikzstyle{arrow} = [thick, ->]
\begin{tikzpicture}[>=latex, node distance=2.5cm]
%     every node/.style = {thin, shape=rectangle, 
%    rounded corners,
%     draw, 
%     align=center},
%    ]   

% \node (geom) [rect, rounded corners] {Caricamento della geometria complessiva del reattore};

\node (compo) [rect, rounded corners] {\textbf{MULTICOMPO}:\\ Database reattoristico che contiene le sezioni d'urto efficaci del combustibile e altri componenti};
\node (ncr) [diam, below of=compo] {modulo \\ NCR};
\node (fmap) [circ, left of=ncr, xshift=-2.5cm] {Fmap};

% \node (ncrd) [rect, rounded corners, below of=ncr] {Interpolazione del database sui parametri di temperatura del combustibile, densità del refrigerante e burn-up};

\node(macrof) [circ, below of=ncr, xshift=-1.5cm] {MacroF};
\node(macro1) [circ, below of=ncr, xshift=1.5cm] {Macro1};

\node(macini) [diam, below of=macrof, xshift=1.5cm] {modulo \\ MACINI};
% \node(macinid) [rect, rounded corners, below of=macini] { byebye };

\node(matex) [circ, left of=macini, xshift=-2.5cm] {Matex};
\node(macro2) [circ, below of=macini] {Macro2};

\node(trivaa) [diam, below of=macro2] {modulo \\ TRIVAA};
\node(system) [circ, below of=trivaa] {System};

\node(track) [circ, left of=trivaa, xshift=-2.5cm] {Track};

\node(flud) [diam, below of=system] {modulo \\ FLUD};
\node(flux) [circ, right of=flud] {Flux};
\node(flpow) [diam, right of=flux] {modulo \\ FLPOW};
\node(power) [circ, right of=flpow] {Power};

\node(rods) [rect, rounded corners, right of=macro1, xshift=2cm] {Parametro per l'inserzione delle barre di controllo: \\ -\textbf{Estratte: rod=0 }\\ -\textbf{Inserite: rod=1}};

\node(para) [rect, rounded corners, right of=ncr, xshift=3.5cm] {Soluzione del campo di temperatura dal codice di termodinamica \\ \textbf{-Temperature del combustibile} \\ \textbf{-Densit\'a del refrigerante}};



\draw [arrow] (compo) -- (ncr);
\draw [arrow] (fmap) -- (ncr);
% \draw [arrow] (ncr) -- (ncrd);
\draw [arrow] (ncr) -- (macrof);
\draw [arrow] (ncr) -- (macro1);
\draw [arrow] (macro1) -- (macini);
\draw [arrow] (macrof) -- (macini);
\draw [arrow] (matex) -- (macini);
% \draw [arrow] (macini) -- (macinid);
\draw [arrow] (macini) -- (macro2);
\draw [arrow] (macro2) -- (trivaa);
\draw [arrow] (track) -- (trivaa);
\draw [arrow] (trivaa) -- (system);
\draw [arrow] (system) -- (flud);
\draw [arrow] (flud) -- (flux);
\draw [arrow] (flux) -- (flpow);
\draw [arrow] (flpow) -- (power);


\draw [arrow] (para) -- (ncr);
\draw [arrow] (rods) -- (macro1);

\end{tikzpicture}
\end{document}
% pdflatex --shell-escape figure.tex
